\documentclass[10pt,a4paper]{amsart}
\usepackage[margin=19mm]{geometry}
\usepackage{amsmath,amssymb,mathtools}
\usepackage[scr=rsfs,cal=euler]{mathalfa}
\usepackage{xfrac}
\usepackage[unicode,hidelinks,bookmarks=false]{hyperref}
\newcommand{\ie}{i.e.\ }
\newcommand{\cf}{cf.\ }
\newcommand{\resp}{respectively\ }
\newcommand{\Subsection}{\subsection}
\providecommand{\nobreakdash}{\nobreak\hbox{-}}
\setlength{\emergencystretch}{2em}
\multlinegap0pt
\usepackage{mathtools}
\usepackage{algorithm}\usepackage{algpseudocode}
\usepackage{enumerate}
\newtheorem{proposition}{Proposition}
\newtheorem{lemma}{Lemma}
\newcommand{\E}{\mathbb{E}}
\newcommand{\R}{\mathbbm{R}}
\providecommand{\mathbbm}{\mathbb} % In case we don't load bbm
\title{Article Title}
\author{Jiaheng Chen, Daniel Sanz-Alonso}
\date{Source: arXiv:2505.24144v2 (CC BY 4.0)}
\begin{document}
\maketitle

\noindent\emph{Source and licence.} Article Title. Version arXiv:2505.24144v2 (CC BY 4.0), distributed by arXiv under the Creative Commons Attribution 4.0 International licence, http://creativecommons.org/licenses/by/4.0/, as recorded in the arXiv licence metadata for this exact version and shown on the abstract page https://arxiv.org/abs/2505.24144v2. Full licence notice: \texttt{licenses/arxiv-2505.24144-CC-BY-4.0.txt}.\par\smallskip

\noindent\textit{Editorial scope.} Counted unit: the article's proposition prop:lowerbound (source0.tex lines 1311--1319). The article's complete original proof of it is delivered with it (source0.tex lines 1321--1436, one \texttt{proof} environment of 116 lines). No mathematics is written, repaired or re-derived: the statement and the proof are the article's own lines, in the article's own notation, and no retained line is substituted or re-flowed.  Retained uncounted as the source-local context the proof needs: lines 2197--2197; lines 2199--2204.  Nothing was omitted from any retained span, so the package carries no omission record.  No retained line contains a citation, so the wrapper carries no bibliography and no bibliography entry.  No retained line contains a figure environment, an image-inclusion command or drawing code, so no image is delivered and the package declares no asset dependency.  Editorial formatting only, in addition: an AMS \texttt{amsart} preamble that restates the article's own declarations and macros used by the retained lines, namely the environments \texttt{proposition}, \texttt{lemma} on separate counters, together with the standard packages those lines need.  The retained environments print in this document as Proposition 1, Lemma 1 in order of appearance, whereas the article prints the same items with the numbering of its own full document.  The complete native source of the article as distributed in the arXiv e-print for this exact version is bundled unchanged as \texttt{sources/179032/source0.tex}.
\section{Proofs of auxiliary results}\label{app:A}

\begin{lemma}\label{lemma:sup_Gaussian_products}
Let $Z^{(k)}_i, 1\le i\le N, 1\le k\le s$ be i.i.d. standard Gaussian random variables in $\R$. Then,
\[
\E  \max_{1\le j\le N} \prod_{k=1}^{s}|Z^{(k)}_j| \asymp_s (\log N)^{s/2}.
\]
\end{lemma}

\begin{proposition}\label{prop:lowerbound} For any integer $p\ge 2$ and $1 \leq k \leq p$, let $X^{(k)}, X_1^{(k)}, \ldots, X_N^{(k)}$ be i.i.d. centered Gaussian random variables in $H^{(k)}$ with covariance operator $\Sigma^{(k)}$, and suppose that $X^{(1)},\ldots, X^{(p)}, (X^{(1)}_i)_{i=1}^{N},\ldots,(X^{(p)}_i)_{i=1}^N,$ are independent. Then,
\[
\E \bigg\|\frac{1}{N} \sum_{i=1}^N X_i^{(1)} \otimes \cdots \otimes X_i^{(p)}-\E\, X^{(1)} \otimes \cdots \otimes X^{(p)}\bigg\| \gtrsim_p\bigg(\prod_{k=1}^p\|\Sigma^{(k)}\|^{1 / 2}\bigg)\mathscr{E}_N\big((\Sigma^{(k)})_{k=1}^p\big),
\]
where
\[
\mathscr{E}_N\big((\Sigma^{(k)})_{k=1}^p\big):= \bigg(\frac{\sum_{k=1}^p r(\Sigma^{(k)})}{N}\bigg)^{1/2}+\frac{1}{N}\prod_{k= 1}^{p}\Big(r(\Sigma^{(k)})+\log N\Big)^{1/2}.
\]
\end{proposition}

\begin{proof}[Proof of Proposition \ref{prop:lowerbound}]
First, we may assume without loss of generality that the effective ranks satisfy $r(\Sigma^{(1)})\ge r(\Sigma^{(2)})\ge \cdots\ge r(\Sigma^{(p)})$. It is straightforward to check that
    \[
\mathscr{E}_N\big((\Sigma^{(k)})_{k=1}^p\big)\asymp_p \begin{cases}
    \Big(\frac{\sum_{k=1}^p r(\Sigma^{(k)})}{N}\Big)^{1/2}+\frac{(\log N)^{(p-t)/2}}{N}\prod_{k=1}^{t}r(\Sigma^{(k)})^{1/2} & \mathrm{ if } \ \exists \ t: r(\Sigma^{(t)})\ge \log N\ge r(\Sigma^{(t+1)}), \\
  \Big(\frac{\sum_{k=1}^p r(\Sigma^{(k)})}{N}\Big)^{1/2}  & \text{otherwise}. \\
\end{cases}
    \]

The conclusion follows directly from the following two claims:

\textbf{Claim I.} It holds that
\begin{align*}
\E \bigg\|\frac{1}{N} \sum_{i=1}^N X_i^{(1)} \otimes \cdots \otimes X_i^{(p)}-\E\, X^{(1)} \otimes \cdots \otimes X^{(p)}\bigg\| \gtrsim_p\bigg(\prod_{k=1}^p\|\Sigma^{(k)}\|^{1 / 2}\bigg) \bigg(\frac{\sum_{k=1}^p r(\Sigma^{(k)})}{N}\bigg)^{1/2}.
\end{align*}

\textbf{Claim II.}
If there exists an index $t$ such that $r(\Sigma^{(t)})\ge \log N \ge r(\Sigma^{(t+1)})$, then
\begin{align*}
\E \bigg\|\frac{1}{N} \sum_{i=1}^N X_i^{(1)} \otimes \cdots \otimes X_i^{(p)}-\E\, X^{(1)} \otimes \cdots \otimes X^{(p)}\bigg\| \gtrsim_p\bigg(\prod_{k=1}^p\|\Sigma^{(k)}\|^{1 / 2}\bigg) \frac{(\log N)^{(p-t)/2}}{N}\prod_{k=1}^{t}r(\Sigma^{(k)})^{1/2}.
\end{align*}

Now we prove \textbf{Claim I} and \textbf{Claim II}.

\textbf{Proof of Claim I:}
By the assumption of independence and the definition of the tensor norm, we obtain
\begin{align*}
&\E \bigg\|\frac{1}{N} \sum_{i=1}^N X_i^{(1)} \otimes \cdots \otimes X_i^{(p)}-\E\, X^{(1)} \otimes \cdots \otimes X^{(p)}\bigg\|=\E \bigg\|\frac{1}{N} \sum_{i=1}^N X_i^{(1)} \otimes \cdots \otimes X_i^{(p)}\bigg\| \\
&= \E \sup_{\|v_1\|= \cdots =\|v_p\|=1} \bigg| \frac{1}{N}\sum_{i=1}^N \langle X^{(1)}_i,v_1\rangle \cdots \langle X^{(p)}_i, v_p \rangle \bigg|\\
&= \E \sup_{\|v_2\|= \cdots =\|v_p\|=1}\bigg\| \frac{1}{N}\sum_{i=1}^N \langle X^{(2)}_i,v_2\rangle \cdots \langle X^{(p)}_i, v_p \rangle X_i^{(1)} \bigg\|\\
&=\E_{(X^{(k)}_i)_{i=1}^{N}, 2\le k\le p} \left[ \E_{(X^{(1)}_i)_{i=1}^{N}}  \sup_{\|v_2\|= \cdots =\|v_p\|=1}\bigg\| \frac{1}{N}\sum_{i=1}^N \langle X^{(2)}_i,v_2\rangle \cdots \langle X^{(p)}_i, v_p \rangle X_i^{(1)} \bigg\|\right].
\end{align*}
Note that, conditionally on $\langle X^{(k)}_i, v_k \rangle, 1\le i\le  N, 2\le k\le p$, the distribution of the random variable
\[
\frac{1}{N}\sum_{i=1}^N \langle X^{(2)}_i,v_2\rangle \cdots \langle X^{(p)}_i, v_p \rangle X_i^{(1)}
\]
is Gaussian and it coincides with the distribution of the random variable
\[
\bigg(\frac{1}{N} \sum_{i=1}^{N}\langle X^{(2)}_i, v_2\rangle^{2}\cdots \langle X^{(p)}_i, v_p\rangle^{2} \bigg)^{1 / 2} \frac{X^{(1)}}{\sqrt{N}}.
\]
Therefore,
\begin{align*}
    &\E_{(X^{(k)}_i)_{i=1}^{N}, 2\le k\le p} \left[ \E_{(X^{(1)}_i)_{i=1}^{N}}  \sup_{\|v_2\|= \cdots =\|v_p\|=1}\bigg\| \frac{1}{N}\sum_{i=1}^N \langle X^{(2)}_i,v_2\rangle \cdots \langle X^{(p)}_i, v_p \rangle X_i^{(1)} \bigg\|\right]\\
    &=\E_{(X^{(k)}_i)_{i=1}^{N}, 2\le k\le p} \left[ \sup_{\|v_2\|= \cdots =\|v_p\|=1}\bigg(\frac{1}{N} \sum_{i=1}^{N}\langle X^{(2)}_i, v_2\rangle^{2}\cdots \langle X^{(p)}_i, v_p\rangle^{2} \bigg)^{1 / 2} \frac{\E \|X^{(1)}\|}{\sqrt{N}}\right]\\
    &=\frac{\E \|X^{(1)}\|}{\sqrt{N}}\E_{(X^{(k)}_i)_{i=1}^{N}, 2\le k\le p}\left[ \sup_{\|v_2\|= \cdots =\|v_p\|=1}\bigg(\frac{1}{N} \sum_{i=1}^{N}\langle X^{(2)}_i, v_2\rangle^{2}\cdots \langle X^{(p)}_i, v_p\rangle^{2} \bigg)^{1 / 2} \right]\\
    &\ge \frac{\E \|X^{(1)}\|}{\sqrt{N}} \sup_{\|v_2\|= \cdots =\|v_p\|=1}\E_{(X^{(k)}_i)_{i=1}^{N}, 2\le k\le p}\bigg(\frac{1}{N} \sum_{i=1}^{N}\langle X^{(2)}_i, v_2\rangle^{2}\cdots \langle X^{(p)}_i, v_p\rangle^{2} \bigg)^{1 / 2}\\
    &=\frac{\E \|X^{(1)}\|}{\sqrt{N}}  \sup_{\|v_2\|= \cdots =\|v_p\|=1} \langle \Sigma^{(2)} v_2,v_2 \rangle^{1/2}\cdots \langle \Sigma^{(p)} v_p,v_p \rangle^{1/2} \E \bigg(\frac{1}{N}\sum_{i=1}^{N} (Z_i^{(2)})^{2} \cdots (Z_i^{(p)})^{2} \bigg)^{1/2}\\
    &=\frac{\E \|X^{(1)}\|}{\sqrt{N}}\bigg(\prod_{k=2}^p\|\Sigma^{(k)}\|^{1 / 2}\bigg)\E \bigg(\frac{1}{N}\sum_{i=1}^{N} (Z_i^{(2)})^{2} \cdots (Z_i^{(p)})^{2} \bigg)^{1/2},
\end{align*}
where
\[
Z^{(k)}_i :=\frac{\langle X^{(k)}_i,v_k \rangle}{\langle \Sigma^{(k)} v_k ,v_k\rangle^{1/2}},\quad 1\le i\le N, \ 1\le k\le p,
\]
are i.i.d. standard Gaussian random variables. One can check that
\[
\E\bigg(\frac{1}{N}\sum_{i=1}^{N} (Z_i^{(2)})^{2} \cdots (Z_i^{(p)})^{2} \bigg)^{1/2}\ge c(p)
\]
for a positive constant $c(p)$ depending only on $p$, which implies that
\begin{align*}
    &\E \bigg\|\frac{1}{N} \sum_{i=1}^N X_i^{(1)} \otimes \cdots \otimes X_i^{(p)}-\E\, X^{(1)} \otimes \cdots \otimes X^{(p)}\bigg\|\\
  %  &=\E_{(X^{(k)}_i)_{i=1}^{N}, 2\le k\le p} \left[ \E_{(X^{(1)}_i)_{i=1}^{N}}  \sup_{\|v_2\|= \cdots =\|v_p\|=1}\bigg\| \frac{1}{N}\sum_{i=1}^N \langle X^{(2)}_i,v_2\rangle \cdots \langle X^{(p)}_i, v_p \rangle X_i^{(1)} \bigg\|\right]\\
    &\ge \frac{\E \|X^{(1)}\|}{\sqrt{N}}\bigg(\prod_{k=2}^p\|\Sigma^{(k)}\|^{1 / 2}\bigg)\E \bigg(\frac{1}{N}\sum_{i=1}^{N} (Z_i^{(2)})^{2} \cdots (Z_i^{(p)})^{2} \bigg)^{1/2}\\
    &\gtrsim_p \bigg(\prod_{k=1}^p\|\Sigma^{(k)}\|^{1 / 2}\bigg)\bigg(\frac{r(\Sigma^{(1)})}{N}\bigg)^{1/2}
    \gtrsim_p \bigg(\prod_{k=1}^p\|\Sigma^{(k)}\|^{1 / 2}\bigg)\bigg(\frac{\sum_{k=1}^p r(\Sigma^{(k)})}{N}\bigg)^{1/2},
\end{align*}
where in the last step we used that $r(\Sigma^{(1)})\ge r(\Sigma^{(2)})\ge \cdots\ge r(\Sigma^{(p)})$. This establishes \textbf{Claim I}.

\textbf{Proof of Claim II:}
Suppose that there exists an index $t$ such that 
\[
r(\Sigma^{(1)})\ge \cdots \ge r(\Sigma^{(t)})\ge \log N \ge r(\Sigma^{(t+1)})\ge \cdots \ge r(\Sigma^{(p)}).
\]
Recall that
\begin{align*}
&\E \bigg\|\frac{1}{N} \sum_{i=1}^N X_i^{(1)} \otimes \cdots \otimes X_i^{(p)}-\E\, X^{(1)} \otimes \cdots \otimes X^{(p)}\bigg\|=\E \bigg\|\frac{1}{N} \sum_{i=1}^N X_i^{(1)} \otimes \cdots \otimes X_i^{(p)}\bigg\| \\
&= \E \sup_{\|v_1\|= \cdots =\|v_p\|=1} \bigg| \frac{1}{N}\sum_{i=1}^N \langle X^{(1)}_i,v_1\rangle \cdots \langle X^{(p)}_i, v_p \rangle \bigg|.
\end{align*}
 

If $t=p$, we have
\begin{align*}
   &\E \sup_{\|v_1\|= \cdots =\|v_p\|=1} \bigg| \frac{1}{N}\sum_{i=1}^N \langle X^{(1)}_i,v_1\rangle \cdots \langle X^{(p)}_i, v_p \rangle \bigg|\\
   &\ge \E_{(X_1^{(k)})_{k=1}^{p}} \sup_{\|v_1\|= \cdots =\|v_p\|=1}  \bigg|\E_{(X_i^{(k)})_{i=2}^{N},1\le k\le p}\frac{1}{N}\sum_{i=1}^N \langle X^{(1)}_i,v_1\rangle \cdots \langle X^{(p)}_i, v_p \rangle \bigg|\\
   &=\E_{(X_1^{(k)})_{k=1}^{p}} \sup_{\|v_1\|= \cdots =\|v_p\|=1}  \bigg| \frac{1}{N} \langle X^{(1)}_1,v_1\rangle \cdots \langle X^{(p)}_1, v_p \rangle \bigg|\\
   &=\frac{1}{N}\E_{(X_1^{(k)})_{k=1}^{p}}\prod_{k=1}^{p}\|X_1^{(k)}\| \overset{\text{($\star$)}}{=} \frac{1}{N}\prod_{k=1}^{p}\E \|X_1^{(k)}\|  \gtrsim_p \frac{1}{N}  \bigg(\prod_{k=1}^p\|\Sigma^{(k)}\|^{1 / 2}\bigg) \prod_{k=1}^{p}r(\Sigma^{(k)})^{1/2},
\end{align*}
as desired. Here $(\star)$ holds since $X^{(1)}_1,\ldots, X^{(p)}_1$ are independent, so $\|X^{(1)}_1\|,\ldots, \|X^{(p)}_1\|$ are independent and the expectation factorizes.


If $1 \le t \le p-1$, then by choosing $v_k=v^{*}_k$ to be the unit eigenvector associated with the largest eigenvalue of $\Sigma^{(k)}$, for each $k \ge t+1$, we obtain 
\begin{align}\label{eq:lowerbound_aux1}
&\E \sup_{\|v_1\|= \cdots =\|v_p\|=1} \bigg| \frac{1}{N}\sum_{i=1}^N \langle X^{(1)}_i,v_1\rangle \cdots \langle X^{(p)}_i, v_p \rangle \bigg|\nonumber\\
    &\ge \E \sup_{\|v_1\|= \cdots =\|v_t\|=1} \bigg| \frac{1}{N}\sum_{i=1}^N \langle X^{(1)}_i,v_1\rangle \cdots \langle X^{(t)}_i, v_t \rangle \langle X^{(t+1)}_i, v^{*}_{t+1} \rangle\cdots \langle X^{(p)}_i, v^{*}_p \rangle\bigg| \nonumber\\
    &= \frac{1}{N}\bigg(\prod_{k=t+1}^p\|\Sigma^{(k)}\|^{1 / 2}\bigg) \E \sup_{\|v_1\|= \cdots =\|v_t\|=1} \bigg| \sum_{i=1}^N \langle X^{(1)}_i,v_1\rangle \cdots \langle X^{(t)}_i, v_t \rangle Z^{(t+1)}_i\cdots Z^{(p)}_i \bigg|,
\end{align}
where \[
Z^{(k)}_i =\frac{\langle X^{(k)}_i,v^*_k \rangle}{\langle \Sigma^{(k)} v^*_k ,v^*_k\rangle^{1/2}}=\frac{\langle X^{(k)}_i,v^*_k \rangle}{\|\Sigma^{(k)}\|^{1/2}},\quad 1\le i\le N, \ 1\le k\le p,
\]
are i.i.d. standard Gaussian random variables.

Observe that
\begin{align}\label{eq:lowerbound_aux2}
    &\E \sup_{\|v_1\|= \cdots =\|v_t\|=1} \bigg| \sum_{i=1}^N \langle X^{(1)}_i,v_1\rangle \cdots \langle X^{(t)}_i, v_t \rangle Z^{(t+1)}_i\cdots Z^{(p)}_i \bigg| \nonumber\\
    &= \E_{(Z^{(k)}_i)_{i=1}^{N},k\ge t+1}\left[\E_{(X^{(k)}_i)_{i=1}^{N},k\le t}  \sup_{\|v_1\|= \cdots =\|v_t\|=1} \bigg| \sum_{i=1}^N \langle X^{(1)}_i,v_1\rangle \cdots \langle X^{(t)}_i, v_t \rangle Z^{(t+1)}_i\cdots Z^{(p)}_i \bigg| \right] \nonumber\\
    &\ge \E_{(Z^{(k)}_i)_{i=1}^{N},k\ge t+1}\left[\max_{1\le j\le N}\E_{(X^{(k)}_j)_{k=1}^{t}} \sup_{\|v_1\|= \cdots =\|v_t\|=1} \bigg|  \E_{(X^{(k)}_i)_{k=1}^{t},i\ne j} \sum_{i=1}^N \langle X^{(1)}_i,v_1\rangle \cdots \langle X^{(t)}_i, v_t \rangle Z^{(t+1)}_i\cdots Z^{(p)}_i \bigg|\right] \nonumber\\
    &=\E_{(Z^{(k)}_i)_{i=1}^{N},k\ge t+1}\left[\max_{1\le j\le N}\E_{(X^{(k)}_j)_{k=1}^{t}} \sup_{\|v_1\|= \cdots =\|v_t\|=1} \bigg| \langle X^{(1)}_j,v_1\rangle \cdots \langle X^{(t)}_j, v_t \rangle Z^{(t+1)}_j\cdots Z^{(p)}_j \bigg|\right] \nonumber\\
    &\ge \E_{(Z^{(k)}_i)_{i=1}^{N},k\ge t+1}\left[\max_{1\le j\le N}\E_{(X^{(k)}_j)_{k=1}^{t}}   \| X^{(1)}_j \| \cdots \| X^{(t)}_j \|
|Z^{(t+1)}_j| \cdots |Z^{(p)}_j| \right] \nonumber\\
&= \E_{(Z^{(k)}_i)_{i=1}^{N},k\ge t+1}\left[\max_{1\le j\le N} (\E \| X^{(1)} \|) \cdots (\E \| X^{(t)} \|)
|Z^{(t+1)}_j| \cdots |Z^{(p)}_j| \right] \nonumber\\
&= \bigg(\prod_{k=1}^{t } \E \| X^{(k)} \|\bigg)\E_{(Z^{(k)}_i)_{i=1}^{N},k\ge t+1} \left[\max_{1\le j\le N} |Z^{(t+1)}_j| \cdots |Z^{(p)}_j| \right] \nonumber\\
&\asymp_p \bigg(\prod_{k=1}^{t } \|\Sigma^{(k)}\|^{1/2}\bigg)\bigg(\prod_{k=1}^{t } r(\Sigma^{(k)})^{1/2}\bigg) (\log N)^{(p-t)/2},
\end{align}
where in the last step we use
Lemma \ref{lemma:sup_Gaussian_products} from Appendix \ref{app:A}, together with the fact that the random variables $Z^{(k)}_i$, for $1\le i\le N$ and $t+1\le k\le p$, are i.i.d. standard Gaussian. Combining \eqref{eq:lowerbound_aux1} and \eqref{eq:lowerbound_aux2} completes the proof of \textbf{Claim II}.   
\end{proof}

\end{document}
